\documentclass[11pt]{article}
\usepackage[margin=1in]{geometry}
\usepackage[T1]{fontenc}
\usepackage[utf8]{inputenc}
\usepackage{listings}
\usepackage{url}
\lstset{breaklines=true,basicstyle=\ttfamily\small,columns=fullflexible}
\setlength{\emergencystretch}{3em}
\title{ES-FA / BTP Local Run Manual}
\author{Commands used for the current repo layout}
\date{\today}
\begin{document}
\maketitle

\section{What This Manual Covers}
This is the short operating manual for the renamed project layout. The working
folders are \path{p1_training/}, \path{p2_hardware_model/}, \path{p3/}, and
\path{output/}. The latest successful smoke run completed the software,
analysis, evidence, and output-generation stages. Hardware/Vivado execution was
intentionally skipped.

\section{Environment Check}
\begin{lstlisting}[language=bash]
.\.venv312\Scripts\python.exe -c "import torch, p1_training.training_core, analysis.compare; print('ok')"
\end{lstlisting}

\section{Start the Whole Local Smoke Pipeline}
\begin{lstlisting}[language=bash]
.\scripts\run_esfa_pipeline.ps1 -PythonExe .\.venv312\Scripts\python.exe -Mode smoke -SkipVivado -SkipHardware
\end{lstlisting}

Successful finish message:
\begin{lstlisting}[language=bash]
ES-FA pipeline complete.
\end{lstlisting}

\section{What the Latest Smoke Run Produced}
\begin{itemize}
\item Best run: \texttt{exp1\_hardware\_aware\_loss}.
\item Accuracy: 95.70\%.
\item Spike sparsity: 56.48\%.
\item Energy proxy improvement vs baseline: -79.68\%.
\item Analysis plots: \texttt{results/plots/}.
\item Evidence tables: \texttt{results/analysis/evidence/}.
\item Paper sections were regenerated in \path{output/}: method, experiment,
tables, key findings, and the draft paper markdown.
\end{itemize}

\section{Manual Software Commands}
\begin{lstlisting}[language=bash]
python p1_training/train_baseline.py
python experiments/run_first2.py
python experiments/run_remaining.py
python iterations/run_all.py
python analysis/compare.py
python analysis/evidence_report.py
python output/generate_paper_output.py
\end{lstlisting}

\section{Hardware Validation Command}
Use this when hardware validation is required:
\begin{lstlisting}[language=bash]
python hardware_validation/kv260/scripts/run_hw_validation.py --model-id baseline_paper1 --scheduler-mode both --clock-mhz 100
\end{lstlisting}

For simulator-only validation:
\begin{lstlisting}[language=bash]
python hardware_validation/kv260/scripts/run_hw_validation.py --model-id baseline_paper1 --scheduler-mode both --clock-mhz 100 --skip-vivado
\end{lstlisting}

\section{CLI Demo}
\begin{lstlisting}[language=bash]
python deployment/cli_interface/main.py --query "hi"
python deployment/cli_interface/main.py --query "strike rate 167.55 balls 39"
python deployment/cli_interface/main.py --query "classify results/runtime/sample_signal.bin"
\end{lstlisting}

\section{Presentation Reminder}
For the current run, present software and analysis results as fresh. Present
hardware/Vivado numbers only as existing local validation evidence, not as a
fresh board run. The KV260 physical measurement stage is still pending.

\end{document}
